\documentclass[12pt]{article}
\usepackage{geometry}
\geometry{
	letterpaper,
	left=20mm,
	right=20mm,
	top=25mm,
	bottom=30mm,
}
\usepackage{graphicx}
\usepackage{tikz}
\usepackage{tikzscale}
\usepackage{pgfplots}

\usepackage{amsmath,amssymb}
\usepackage{enumitem}
\usepackage{algorithm}
\usepackage{algorithmicx}
\usepackage{algpseudocode}
\usepackage{subcaption}
\usepackage{mathtools}
\usepackage{amsthm}
\usepackage{nccmath}
\usepackage{listings}
\usepackage{tcolorbox}
\usepackage{booktabs}
\usepackage{tabularx}
\usepackage{longtable}
\usepackage{multicol}

\makeatletter
\def\verbatim@font{\linespread{1}\normalfont\ttfamily}
\makeatother
\usepackage[breaklinks]{hyperref}
\hypersetup{
	colorlinks = true,
	citecolor = blue,
	linkcolor = blue
}
\usepackage[T1]{fontenc}
\DeclarePairedDelimiter{\ceil}{\lceil}{\rceil}
\mathchardef\mhyphen="2D
\newcommand\makeset{\mathop{make\mhyphen set}}
\newcommand\findset{\mathop{find\mhyphen set}}
\newcommand*{\Comb}[2]{{}^{#1}C_{#2}}
\usepackage{mdframed}
\usepackage{fancyhdr}
\fancyhf{}
\chead{\courseName}
\lhead{\thepage}
\rhead{\homeworkName}
\renewcommand{\headrulewidth}{1pt}
\renewcommand{\footrulewidth}{0pt}
\fancypagestyle{first}{
	\fancyhf{}
	\chead{\small}
	\renewcommand{\headrulewidth}{0pt}
	\renewcommand{\footrulewidth}{0pt}
}
\fancypagestyle{second}{
	\fancyhf{}
	\chead{}
	\renewcommand{\headrulewidth}{0pt}
	\renewcommand{\footrulewidth}{0pt}
}
\pagestyle{fancy}

\newcommand{\mBra}[1]{\vert #1 \rangle}
\newcommand{\mKet}[1]{\langle #1 \vert}
\newcommand{\mBraKet}[2]{\langle #1 \vert #2 \rangle}
\newcommand{\mBraXKet}[3]{\langle #1 \vert #2 \vert #3 \rangle}
\newcommand{\qOne}{\mBra{1}}
\newcommand{\qZero}{\mBra{0}}
\newcommand{\qPlus}{\mBra{+}}
\newcommand{\qMinus}{\mBra{-}}
\newcommand{\qOneKet}{\mKet{1}}
\newcommand{\qZeroKet}{\mKet{0}}
\newcommand{\qPlusKet}{\mKet{+}}
\newcommand{\qMinusKet}{\mKet{-}}
\newcommand{\mAbs}[1]{\vert #1 \vert}
\newcommand{\mAbsSq}[1]{\mAbs{#1}^{2}}
\newcommand{\mSecNorm}[1]{\vert\vert #1 \vert\vert}
\newcommand{\iu}{{i\mkern1mu}}


\usepackage{xepersian}
\settextfont[
    Scale=1.2,
    Path={LaTeX_Common/fonts/BNazanin/},
    Extension=.ttf,
    BoldFont = BNAZANB
]{BNAZANIN}
\setlatintextfont[
    Scale=1.1,
    Path={LaTeX_Common/fonts/LatinModernRoman/},
    Extension=.otf
]{LMR_Regular}

\newcommand{\courseName}{شبکه‌های کامپیوتری پیشرفته}
\newcommand{\courseSemester}{پاییز ۱۴۰۵}
\newcommand{\homeworkName}{راهنمای درس}
\newcommand{\homeworkDue}{دانشکده مهندسی کامپیوتر}

\renewcommand{\baselinestretch}{1.5}

\newcommand{\en}[1]{\lr{#1}}
\newcommand{\pending}{\textbf{بعداً اعلام می‌شود}}
\newcommand{\projectarea}[4]{%
\begin{tcolorbox}[
    colback=black!2,
    colframe=black!45,
    boxrule=0.55pt,
    arc=1.2mm,
    left=2mm,right=2mm,top=1.5mm,bottom=1.5mm
]
{\large\bfseries\lr{#1}}\par
\smallskip
\textbf{ده جنبه‌ی پرتکرار طراحی}\par
\begin{multicols}{2}
\begin{enumerate}[leftmargin=*,nosep,label=\arabic*.]
#2
\end{enumerate}
\end{multicols}
\smallskip
\textbf{برچسب‌های فناوری و بافت:} \lr{#3}
\par\smallskip
\textbf{مقاله‌های کلیدی برای شروع:}\par
\begin{latin}
\small
\begin{itemize}[leftmargin=*,nosep]
#4
\end{itemize}
\end{latin}
\end{tcolorbox}
\vspace{0.5em}
}

\begin{document}
	\graphicspath{{LaTeX_Common/cover/}}
	\pagenumbering{gobble}
	\thispagestyle{first}
\begin{mdframed}
\begin{minipage}[t]{0.15\textwidth}
	\centering
	\includegraphics[width=0.7\textwidth]{sharif} \\
	\homeworkName
\end{minipage}%
\begin{minipage}[b]{0.64\textwidth}
	\centering
	\courseName \\
	\courseSemester
\end{minipage}%
\hfill
\begin{minipage}[t]{0.15\textwidth}
	\centering
	\includegraphics[width=0.7\textwidth]{sharif} \\
	\homeworkDue
\end{minipage}%
\end{mdframed}


	\pagenumbering{arabic}

\section{هدف و ساختار کلی درس}
این درس بر تصمیم‌های طراحی در شبکه‌ها و سامانه‌های شبکه‌ای پیشرفته تمرکز دارد. هدف فقط آشنایی با مجموعه‌ای از فناوری‌ها نیست؛ انتظار می‌رود بتوانید درباره‌ی انتخاب‌های طراحی، مقیاس‌پذیری، قابلیت اطمینان، کارایی، روش ارزیابی و محدودیت‌های یک سامانه استدلال کنید و نتیجه‌ی خود را با شواهد دفاع‌پذیر پشتیبانی کنید.

ارزیابی درس از چهار جزء اصلی تشکیل می‌شود:
\begin{itemize}[label=$\bullet$]
    \item مجموعه‌ای از مقاله‌های پژوهشی مشترک کلاس؛ \textbf{چهار مقاله‌ی پژوهشی مشترک کلاس (شماره‌های \lr{\#2} تا \lr{\#5})}؛
    \item یک \textbf{پروژه‌ی گروهی} با ساختاری نزدیک به یک پروژه‌ی پژوهشی؛
    \item یک \textbf{آزمون میان‌ترم}؛
    \item یک \textbf{آزمون پایان‌ترم}.
\end{itemize}

نسبت نمره‌ی این اجزا در نسخه‌ی فعلی این راهنما نهایی نشده و \pending. زمان‌بندی مقاله‌های مشترک، جلسه‌های پروژه، میان‌ترم و پایان‌ترم در جدول زیر آمده است. بازه‌ی دقیق جلسه‌ی هر گروه با دستیار آموزشی هماهنگ خواهد شد.

\section{تقویم درس و موعدهای اصلی}
این تقویم بر اساس برنامه‌ی دستیاران آموزشی درس تنظیم شده است. چهار مقاله‌ی مشترک کلاس به‌ترتیب با شماره‌های \lr{\#2}، \lr{\#3}، \lr{\#4} و \lr{\#5} مشخص می‌شوند؛ \textbf{مقاله‌ی \lr{\#1} در برنامه‌ی فعلی وجود ندارد.} زمان‌های جلسه‌های پروژه بازه هستند و زمان دقیق هر گروه را دستیار آموزشی هماهنگ می‌کند.

\begingroup
\small
\renewcommand{\arraystretch}{1.18}
\begin{longtable}{>{\raggedleft\arraybackslash}p{0.52\textwidth} >{\raggedleft\arraybackslash}p{0.19\textwidth} >{\raggedleft\arraybackslash}p{0.20\textwidth}}
\toprule
\textbf{رویداد} & \textbf{تاریخ شمسی ۱۴۰۵} & \textbf{تاریخ میلادی}\\
\midrule
\endfirsthead
\toprule
\textbf{رویداد} & \textbf{تاریخ شمسی ۱۴۰۵} & \textbf{تاریخ میلادی}\\
\midrule
\endhead
اعلام و انتخاب پروژه‌ها & ۹ آبان & \lr{31 Oct 2026}\\
اعلام مقاله‌ی \lr{\#2} & ۱۱ آبان & \lr{2 Nov 2026}\\
بازه‌ی جلسه‌ی اول پروژه و آغاز فاز اول & ۱۸ تا ۲۲ آبان & \lr{9--13 Nov 2026}\\
موعد مقاله‌ی \lr{\#2} & ۲۵ آبان & \lr{16 Nov 2026}\\
اعلام مقاله‌ی \lr{\#3} & ۳ آذر & \lr{24 Nov 2026}\\
بازه‌ی جلسه‌ی دوم پروژه و آغاز فاز دوم & ۷ تا ۹ آذر & \lr{28--30 Nov 2026}\\
\textbf{آزمون میان‌ترم} & \textbf{۱۴ آذر} & \lr{\textbf{5 Dec 2026}}\\
اعلام مقاله‌ی \lr{\#4} & ۱۶ آذر & \lr{7 Dec 2026}\\
موعد مقاله‌ی \lr{\#3} & ۱۷ آذر & \lr{8 Dec 2026}\\
بازه‌ی جلسه‌ی سوم پروژه و آغاز فاز سوم & ۲۸ تا ۳۰ آذر & \lr{19--21 Dec 2026}\\
بازه‌ی جلسه‌ی چهارم و جمع‌بندی پروژه & ۱۴ تا ۱۸ دی & \lr{4--8 Jan 2027}\\
موعد مقاله‌ی \lr{\#4} و اعلام مقاله‌ی \lr{\#5} & ۲۱ دی & \lr{11 Jan 2027}\\
\textbf{آزمون پایان‌ترم} & \textbf{۲۶ دی} & \lr{\textbf{16 Jan 2027}}\\
موعد مقاله‌ی \lr{\#5} & ۱۲ بهمن & \lr{1 Feb 2027}\\
\bottomrule
\end{longtable}
\endgroup

\noindent\textbf{توجه:} موعد مقاله‌ی \lr{\#5} بعد از آزمون پایان‌ترم است. اگر زمان یکی از بازه‌های جلسه‌ی پروژه یا تحویل‌ها تغییر کند، اطلاع‌رسانی رسمی درس ملاک خواهد بود.

\section{مطالعه‌ی مقاله‌های مشترک کلاس}
در طول ترم \textbf{چهار مقاله‌ی مشترک با شماره‌های \lr{\#2} تا \lr{\#5}} تعیین می‌شود. هدف از مطالعه صرفاً خلاصه‌کردن متن مقاله نیست. انتظار می‌رود بتوانید مسئله، انتخاب‌های طراحی، روش ارزیابی، شواهد اصلی، محدودیت‌ها و رابطه‌ی آن با مفاهیم درس را توضیح دهید.

این چهار مقاله از مقاله‌های پژوهشی‌ای که برای پروژه‌ی گروهی انتخاب می‌کنید جدا هستند. زمان اعلام و موعد هر مقاله در جدول بالا مشخص است؛ عنوان و دستورالعمل ارزیابی هر مقاله جداگانه ارائه خواهد شد.

\section{پروژه‌ی گروهی}
\subsection{ساختار تیم}
پروژه در گروه‌های \textbf{سه یا چهار نفره} انجام می‌شود. هر گروه یک حوزه‌ی کلی را انتخاب می‌کند و سپس خودش مشخص می‌کند کدام دامنه‌ها، جنبه‌های طراحی و فناوری‌ها را در آن حوزه بررسی خواهد کرد.

اگر یکی از اعضای گروه در طول ترم از درس حذف شود، انتظار نمی‌رود اعضای باقی‌مانده به‌صورت خودکار مسئولیت کامل او را نیز بر عهده بگیرند. دامنه و برنامه‌ی پروژه با کمک دستیار آموزشی بازتنظیم می‌شود تا کار مربوط به مسئولیت دانشجوی حذف‌شده محدود یا حذف شود و حجم باقی‌مانده متناسب با اندازه‌ی جدید گروه باشد.

\subsection{ماهیت پروژه}
پروژه باید از نظر شیوه‌ی کار به یک پروژه‌ی پژوهشی نزدیک باشد. گروه یک پاسخ ازپیش‌تعیین‌شده یا فهرست کامل وظایف دریافت نمی‌کند. دانشجویان باید مسئله را صورت‌بندی کنند، فضای طراحی را بشناسند، پرسش پژوهشی بسازند، گزینه‌های قابل بررسی را انتخاب کنند، آزمایش یا تحلیل مناسب طراحی کنند و در طول پروژه بر اساس شواهد مسیر خود را اصلاح کنند.

واحد اصلی برای تعریف پروژه چنین است:
\begin{center}
\textbf{دامنه $\times$ جنبه‌ی طراحی $\times$ برچسب‌های اختیاری فناوری و بافت}
\end{center}

گروه بر اساس این ابعاد یک \textbf{ماتریس پروژه} می‌سازد. ماتریس قرار نیست همه‌ی خانه‌های ممکن را پوشش دهد؛ هدف آن است که فضای مسئله را روشن کند و نشان دهد گروه آگاهانه چه بخش‌هایی را برای بررسی انتخاب کرده است.

\subsubsection*{نمونه‌ی ماتریس پروژه}
برای نمونه، فرض کنید یک گروه حوزه‌ی \lr{Real-Time Media \& Immersive Networking} را انتخاب کرده است. یک ماتریس اولیه می‌تواند چند زیرحوزه را در ستون‌ها و چند جنبه‌ی طراحی را در سطرها قرار دهد. خانه‌های «اصلی» مستقیماً در پرسش و آزمایش‌های پروژه بررسی می‌شوند، خانه‌های «مقایسه» فقط برای ایجاد خط مبنا یا درک بهتر استفاده می‌شوند، و خانه‌ی خالی خارج از دامنه‌ی فعلی پروژه است.

\begin{center}
\small
\renewcommand{\arraystretch}{1.35}
\begin{tabularx}{\textwidth}{
    >{\raggedleft\arraybackslash}p{0.28\textwidth}
    >{\centering\arraybackslash}X
    >{\centering\arraybackslash}X
    >{\centering\arraybackslash}X
}
\toprule
\textbf{جنبه‌ی طراحی} &
\textbf{\lr{Video Conferencing}} &
\textbf{\lr{Cloud Gaming}} &
\textbf{\lr{XR / AR / VR}} \\
\midrule
\lr{Latency \& Deadline Optimization} &
مقایسه & \textbf{اصلی} & \textbf{اصلی} \\
\lr{Transport \& Congestion Control} &
{} & \textbf{اصلی} & \textbf{اصلی} \\
\lr{Application Adaptation \& Cross-Layer Control} &
مقایسه & مقایسه & \textbf{اصلی} \\
\lr{Placement, Offloading \& Partitioning} &
{} & مقایسه & \textbf{اصلی} \\
\bottomrule
\end{tabularx}
\end{center}

برای چنین ماتریسی، یک پرسش اولیه می‌تواند این باشد که در یک سرویس تعاملی \lr{XR} با بخشی از پردازش در لبه، ترکیب کنترل ازدحام، سازگارسازی کیفیت و محل پردازش تحت بارهای مختلف شبکه چه اثری بر تأخیر انتهابه‌انتها و کیفیت تجربه دارد. این فقط یک نمونه است؛ ماتریس هر گروه باید متناسب با حوزه، پرسش و روش ارزیابی خودش ساخته شود و لازم نیست شکل یا ابعاد یکسانی داشته باشد.

\subsection{مسئولیت فردی درون پروژه}
هر عضو گروه باید مسئولیت مشخصی در بخشی از دامنه یا یکی از جنبه‌های طراحی پروژه داشته باشد. این مسئولیت نباید به معنای کارکردن کاملاً مستقل از گروه باشد؛ هر دانشجو باید بداند کار خودش چگونه به پرسش پژوهشی و نتیجه‌ی کلی پروژه متصل می‌شود.

در ارزیابی پروژه، علاوه بر خروجی تیم، به دو موضوع فردی توجه می‌شود:
\begin{itemize}[label=$\bullet$]
    \item میزان تسلط دانشجو بر دامنه یا جنبه‌ی طراحی‌ای که مسئول آن است و توانایی توضیح اثر آن بر کل پروژه؛
    \item کیفیت انجام وظایفی که در هر مرحله به او واگذار شده است، بر اساس تعریف انجام و شواهد توافق‌شده برای همان وظیفه.
\end{itemize}

\section{دستیار آموزشی و چرخه‌ی پروژه}
برای هر پروژه یک دستیار آموزشی با عنوان \en{TA} تعیین می‌شود که در طول پروژه همان گروه را راهنمایی می‌کند. نقش دستیار آموزشی انتخاب مسیر به جای دانشجویان نیست؛ گروه باید یافته‌ها، تصمیم‌های پیشنهادی و ایده‌ی ادامه‌ی کار را به جلسه بیاورد. دستیار آموزشی به اصلاح دامنه، تبدیل ایده‌ها به وظایف قابل ارزیابی و حفظ تناسب حجم کار کمک می‌کند.

در مجموع \textbf{چهار جلسه‌ی پروژه با دستیار آموزشی} برگزار می‌شود. میان این چهار جلسه، \textbf{سه فاز کاری پروژه} قرار دارد:
\begin{center}
جلسه‌ی ۱ $\longrightarrow$ فاز ۱ $\longrightarrow$ جلسه‌ی ۲ $\longrightarrow$ فاز ۲ $\longrightarrow$ جلسه‌ی ۳ $\longrightarrow$ فاز ۳ $\longrightarrow$ جلسه‌ی ۴
\end{center}

همه‌ی اعضای تیم باید در این جلسات حضور داشته باشند. در هر جلسه:
\begin{enumerate}[label=\arabic*.]
    \item گروه یافته‌ها و پیشرفت از جلسه‌ی قبلی را ارائه می‌کند؛
    \item هر دانشجو وضعیت دامنه یا جنبه‌ی طراحیِ تحت مسئولیت خود، وظایف شخصی انجام‌شده و معنای نتایج آن‌ها برای پروژه را توضیح می‌دهد؛
    \item گروه توضیح می‌دهد بر اساس شواهد فعلی چگونه می‌خواهد پروژه را ادامه دهد؛
    \item مقاله‌های پژوهشی مرتبط و اثری که بر انتخاب مسیر بعدی داشته‌اند بررسی می‌شوند؛
    \item دستیار آموزشی درباره‌ی دامنه، اعتبارسنجی و حجم کار بازخورد می‌دهد و در صورت بیش‌ازحد بزرگ یا کوچک بودن محدوده، درخواست اصلاح می‌کند؛
    \item برای فاز بعد، وظایف مشخص با مسئول فردی و تعریف انجام روشن ساخته و ثبت می‌شوند.
\end{enumerate}

وظایف باید به پرسش پروژه متصل باشند و خروجی آن‌ها قابل بررسی باشد. یک وظیفه ممکن است پیاده‌سازی، آزمایش، بازتولید یک نتیجه، اندازه‌گیری، مقایسه، تحلیل، ساخت ابزار، بررسی یک فرضیه یا اعتبارسنجی یک تصمیم طراحی باشد. نتیجه‌ی منفی نیز می‌تواند یک نتیجه‌ی معتبر باشد، اگر روش درست اجرا و تحلیل شده باشد.

\section{جلسه‌ی اول پروژه}
جلسه‌ی اول نقطه‌ی شروع رسمی پروژه است. گروه باید پیش از این جلسه حداقل موارد زیر را آماده کرده باشد:
\begin{itemize}[label=$\square$]
    \item یک یا چند \textbf{پرسش پژوهشی اولیه}؛
    \item \textbf{ماتریس پروژه} شامل دامنه‌ها، جنبه‌های طراحی و در صورت نیاز برچسب‌های فناوری؛
    \item مشخص‌کردن خانه‌ها یا ترکیب‌هایی که گروه قصد دارد بررسی کند؛
    \item توضیح اولیه‌ی این‌که گروه چگونه می‌خواهد پرسش را بررسی کند: ابزار، محیط آزمایش، داده، خط مبنا، روش اندازه‌گیری یا نوع تحلیل؛
    \item تقسیم اولیه‌ی مسئولیت‌های دامنه یا جنبه‌ی طراحی میان اعضای تیم.
\end{itemize}

در خود جلسه، این طرح با دستیار آموزشی بررسی و در صورت نیاز کوچک یا بازساخت می‌شود. سپس وظایف فاز اول با کمک دستیار آموزشی تعریف می‌شوند و بر اساس مسئولیت هر دانشجو به اعضای تیم واگذار می‌شوند.

\section{مقاله‌های پروژه}
در هر یک از سه فاز پروژه، \textbf{هر دانشجو باید حداقل یک مقاله‌ی پژوهشی مرتبط با پروژه} را تا جلسه‌ی بعدی مطالعه کرده باشد؛ در صورت مفیدبودن می‌تواند مقاله را زودتر از موعد مطرح کند. در نتیجه، هر دانشجو در طول سه فاز حداقل سه مقاله‌ی پروژه‌محور بررسی می‌کند.

هدف این مقاله‌ها ایجاد ورودی مستقیم برای تصمیم‌های پروژه است. دانشجو باید بتواند توضیح دهد:
\begin{itemize}[label=$\bullet$]
    \item مقاله چه مسئله‌ای را حل می‌کند و کدام بخش آن به پروژه‌ی گروه مرتبط است؛
    \item چه فرض، طراحی، روش، خط مبنا، معیار یا نتیجه‌ای از مقاله برای پروژه مفید است؛
    \item مقاله چه تغییری در ماتریس، فرضیه، روش آزمایش یا وظایف پیشنهادی بعدی ایجاد می‌کند؛
    \item محدودیت‌های مقاله چیست و کدام بخش آن نباید بدون اعتبارسنجی مستقیماً پذیرفته شود.
\end{itemize}

فهم دانشجو از مقاله در جلسه توسط دستیار آموزشی بررسی و به‌عنوان بخشی از ارزیابی پروژه ثبت می‌شود. بنابراین، صرف ارائه‌ی خلاصه یا نام مقاله کافی نیست.

\section{استفاده از هوش مصنوعی در پروژه}
ابزارهای هوش مصنوعی و عامل‌های کدنویسی می‌توانند برای پیاده‌سازی، جست‌وجوی اولیه، تحلیل و توسعه‌ی آزمایش‌ها استفاده شوند. مسئولیت نهایی فهم، صحت‌سنجی و دفاع از خروجی همچنان با دانشجویان است. هر نتیجه‌ای که وارد گزارش یا استدلال پروژه می‌شود باید توسط تیم قابل توضیح و اعتبارسنجی باشد.

جلسات پروژه به‌صورت گروهی برگزار می‌شوند و هر عضو باید بتواند درباره‌ی تصمیم‌ها و شواهد مربوط به مسئولیت خودش پاسخ دهد؛ استفاده از ابزار هوش مصنوعی جایگزین این مسئولیت فردی نیست.

\section{خروجی و ارزیابی پروژه}
جزئیات وزن نمره‌ی پروژه در نمره‌ی کل درس همراه با ساختار نهایی ارزیابی اعلام می‌شود. در خود پروژه، ارزیابی بر کیفیت پژوهش و فهم دانشجو متمرکز است، نه صرفاً بر حجم کد یا تعداد وظایف.

موارد مهم عبارت‌اند از:
\begin{itemize}[label=$\bullet$]
    \item کیفیت صورت‌بندی پرسش و انتخاب دامنه‌ی شدنی؛
    \item توانایی استفاده از ماتریس برای توضیح فضای طراحی و انتخاب‌های بررسی‌شده؛
    \item کیفیت وظایف تعریف‌شده، انجام آن‌ها و شواهد ارائه‌شده؛
    \item تسلط هر دانشجو بر مسئولیت دامنه یا جنبه‌ی طراحی خود؛
    \item استفاده‌ی معنادار از مقاله‌های پروژه برای تصمیم‌گیری؛
    \item کیفیت روش ارزیابی، خطوط مبنا، اندازه‌گیری‌ها و اعتبارسنجی؛
    \item توانایی گروه در استنتاج نتیجه از شواهد و بیان محدودیت‌ها و عدم‌قطعیت‌ها.
\end{itemize}

\appendix
\section{حوزه‌های پیشنهادی پروژه، جنبه‌های طراحی و برچسب‌های فناوری}
اطلاعات لازم برای انتخاب اولیه‌ی حوزه‌ی پروژه در همین پیوست آمده است. هر گروه ابتدا یکی از \textbf{سیزده حوزه‌ی مجاز} را انتخاب می‌کند، سپس با استفاده از جنبه‌های طراحی و برچسب‌های فناوری، محدوده‌ی دقیق‌تری برای ماتریس و پرسش پژوهشی خود می‌سازد. حوزه‌های \lr{Cellular \& Wireless Access}، \lr{Satellite \& Non-Terrestrial Networks} و \lr{Harsh-Environment \& Delay-Tolerant Networks} در درس فعلی \textbf{خارج از دامنه} هستند. مقاله‌های ذکرشده در هر حوزه، منابع اولیه‌ی منتخب برای شروع مطالعه‌اند و به‌خودی‌خود موضوع یا محدوده‌ی نهایی پروژه را تعیین نمی‌کنند.

فهرست «ده جنبه‌ی پرتکرار طراحی» برای کمک به شناخت سریع فضای پژوهش هر حوزه مرتب شده است. ترتیب آن‌ها فقط نشان‌دهنده‌ی فراوانی نسبی موضوع در مرور پژوهش‌های اخیر است و به معنی بهتر بودن، آسان‌تر بودن یا الزامی بودن آن جنبه نیست. برچسب‌های فناوری نیز برای محدودکردن دامنه یا ساختن یک بافت مشخص برای پروژه پیشنهاد شده‌اند.

\begin{tcolorbox}[
    colback=black!3,
    colframe=black!50,
    boxrule=0.6pt,
    arc=1.2mm,
    left=2.5mm,right=2.5mm,top=2mm,bottom=2mm
]
\textbf{روش استفاده از این پیوست:}
یک یا چند خانه‌ی معنادار از ترکیب \textbf{حوزه $\times$ جنبه‌ی طراحی} را انتخاب کنید و در صورت نیاز با یک یا چند برچسب فناوری محدودترشان کنید. لازم نیست همه‌ی جنبه‌ها را پوشش دهید. اگر چند جنبه را هم‌زمان انتخاب می‌کنید، باید تعامل میان آن‌ها بخشی واقعی از پرسش پروژه باشد و کل محدوده همچنان در سه فاز پروژه قابل بررسی بماند.
\end{tcolorbox}



\projectarea{Edge, Cloud \& Networked Services}{
    \item \lr{Placement, Offloading \& Partitioning}
    \item \lr{Scheduling \& Resource Allocation}
    \item \lr{Security, Privacy \& Trust}
    \item \lr{Mobility \& Handover}
    \item \lr{Latency \& Deadline Optimization}
    \item \lr{Caching \& Content Placement}
    \item \lr{MAC, Spectrum \& Radio Resource Management}
    \item \lr{Cost, Economics \& Provisioning}
    \item \lr{Coding, Compression \& Redundancy}
    \item \lr{Energy, Carbon \& Sustainability}
}{MEC, edge intelligence, cloud-edge orchestration, fog computing, serverless, microservices, service mesh, digital twins, disaggregated resources}{
\item Autothrottle: A Practical Bi-Level Approach to Resource Management for SLO-Targeted Microservices (\textit{NSDI 2024})
\item Microns: Connection Subsetting for Microservices in Shared Clusters (\textit{SIGMETRICS 2025})
}

\projectarea{Packet Networks \& End-to-End Protocols}{
    \item \lr{Transport \& Congestion Control}
    \item \lr{Routing \& Path Selection}
    \item \lr{Security, Privacy \& Trust}
    \item \lr{Scheduling \& Resource Allocation}
    \item \lr{Reliability, Fault Tolerance \& Recovery}
    \item \lr{Protocol Design \& Interoperability}
    \item \lr{Traffic Engineering \& Load Balancing}
    \item \lr{Latency \& Deadline Optimization}
    \item \lr{Cost, Economics \& Provisioning}
    \item \lr{Coding, Compression \& Redundancy}
}{transport protocols, packet routing, queueing, multipath, TCP, QUIC, FEC, SDN-assisted control}{
\item No Signal to Rule Them All: A Systematic Analysis of In-Network Congestion Signals (\textit{NINeS 2026})
\item Making Congestion Control Algorithms Insensitive to Underlying Propagation Delays (\textit{NINeS 2026})
}

\projectarea{Autonomous \& Cyber-Physical Networks}{
    \item \lr{Scheduling \& Resource Allocation}
    \item \lr{Placement, Offloading \& Partitioning}
    \item \lr{Security, Privacy \& Trust}
    \item \lr{Mobility \& Handover}
    \item \lr{Routing \& Path Selection}
    \item \lr{MAC, Spectrum \& Radio Resource Management}
    \item \lr{Latency \& Deadline Optimization}
    \item \lr{Caching \& Content Placement}
    \item \lr{Energy, Carbon \& Sustainability}
    \item \lr{Architecture \& System Design}
}{vehicular/V2X, UAV/aerial, robotics/teleoperation, V2X, UAV swarms, autonomous driving, cooperative perception, ITS, digital twins, MEC, 5G/6G}{
\item CoPe: Taming Collaborative 3D Perception via Lite Network Attention across Mobile Agents (\textit{ICDCS 2025})
\item TURBO: Utility-Aware Bandwidth Allocation for Cloud-Augmented Autonomous Control (\textit{NINeS 2026})
}
\noindent{\footnotesize\textbf{شرط انتخاب:} این حوزه فقط زمانی پذیرفته است که مسئله‌ی اصلی و معیارهای ارزیابی بر \textbf{شبکه و تصمیم‌های طراحی شبکه} متمرکز باشند؛ پروژه‌های صرفاً مربوط به رباتیک، ادراک یا کنترل خودران خارج از دامنه‌اند.}\par


\projectarea{IoT, Sensor \& Industrial Networks}{
    \item \lr{Scheduling \& Resource Allocation}
    \item \lr{Security, Privacy \& Trust}
    \item \lr{Routing \& Path Selection}
    \item \lr{Placement, Offloading \& Partitioning}
    \item \lr{Latency \& Deadline Optimization}
    \item \lr{MAC, Spectrum \& Radio Resource Management}
    \item \lr{Localization \& Positioning}
    \item \lr{Throughput, Capacity \& Scalability}
    \item \lr{Energy, Carbon \& Sustainability}
    \item \lr{Topology, Fabric \& Architecture}
}{sensor networks, LPWAN, industrial/factory networks, deterministic networking, TSN, LoRa/LPWAN, RFID, backscatter, energy harvesting, industrial IoT, URLLC}{
\item T(SN)-Ray: Gauging TAS and PSFP Delays of TSN Switches for Predictable Deterministic Networking (\textit{CoNEXT 2025})
\item Elastic Scheduling for Mix-Flow in Time-Sensitive Networking (\textit{ICDCS 2025})
}

\projectarea{Datacenter \& AI Cluster Networks}{
    \item \lr{Transport \& Congestion Control}
    \item \lr{Scheduling \& Resource Allocation}
    \item \lr{Programmability \& Hardware Offload}
    \item \lr{Placement, Offloading \& Partitioning}
    \item \lr{Routing \& Path Selection}
    \item \lr{Traffic Engineering \& Load Balancing}
    \item \lr{Latency \& Deadline Optimization}
    \item \lr{Architecture \& System Design}
    \item \lr{Coding, Compression \& Redundancy}
    \item \lr{Topology, Fabric \& Architecture}
}{datacenter transport, AI/GPU clusters, RDMA, RoCE, ECN, PFC, SmartNIC/DPU, collectives, coflows, in-network aggregation, optical interconnects, P4}{
\item Hawkeye: Diagnosing RDMA Network Performance Anomalies with PFC Provenance (\textit{SIGCOMM 2025})
\item SyCCL: Exploiting Symmetry for Efficient Collective Communication Scheduling (\textit{SIGCOMM 2025})
}

\projectarea{Network Security, Privacy \& Censorship}{
    \item \lr{Security, Privacy \& Trust}
    \item \lr{Traffic Analysis \& Classification}
    \item \lr{Control-Plane Design \& Management}
    \item \lr{Programmability \& Hardware Offload}
    \item \lr{Scheduling \& Resource Allocation}
    \item \lr{Telemetry \& Monitoring}
    \item \lr{Architecture \& System Design}
    \item \lr{Coding, Compression \& Redundancy}
    \item \lr{Measurement \& Characterization}
    \item \lr{Packet Processing \& Data Structures}
}{routing/DNS security, traffic security, censorship resistance, anonymity, DDoS, intrusion detection, encrypted traffic, DNSSEC, TLS, QUIC, Tor/onion services}{
\item A System to Detect Forged-Origin BGP Hijacks (\textit{NSDI 2024})
\item Exposing and Circumventing SNI-based QUIC Censorship of the Great Firewall of China (\textit{USENIX Security 2025})
}

\projectarea{Real-Time Media \& Immersive Networking}{
    \item \lr{Latency \& Deadline Optimization}
    \item \lr{Architecture \& System Design}
    \item \lr{Application Adaptation \& Cross-Layer Control}
    \item \lr{Transport \& Congestion Control}
    \item \lr{QoS, QoE \& SLO Management}
    \item \lr{Scheduling \& Resource Allocation}
    \item \lr{Coding, Compression \& Redundancy}
    \item \lr{Placement, Offloading \& Partitioning}
    \item \lr{Measurement \& Characterization}
    \item \lr{Caching \& Content Placement}
}{conferencing, interactive/live streaming, cloud gaming, XR/AR/VR, volumetric media, WebRTC, ABR, FEC, QUIC, volumetric video, 3D Gaussian splatting}{
\item Syntra: Synthesizing Cross-Layer Controllers for Low-Latency Video Streaming (\textit{NSDI 2026})
\item EdgeGaussian: Real-time Free-Viewpoint Video for Mobile VR via Edge-Client Collaborative Neural Rendering (\textit{MobiCom 2025})
}

\projectarea{Distributed AI \& LLM Systems}{
    \item \lr{Security, Privacy \& Trust}
    \item \lr{Scheduling \& Resource Allocation}
    \item \lr{Coding, Compression \& Redundancy}
    \item \lr{Model/State Placement \& Movement}
    \item \lr{Placement, Offloading \& Partitioning}
    \item \lr{Cost, Economics \& Provisioning}
    \item \lr{Latency \& Deadline Optimization}
    \item \lr{Energy, Carbon \& Sustainability}
    \item \lr{Caching \& Content Placement}
    \item \lr{Throughput, Capacity \& Scalability}
}{distributed/federated training, inference/serving, split inference, federated learning, distributed training, LLM inference, model serving, MoE, KV cache, collectives, differential privacy}{
\item MegaScale-Infer: Efficient Mixture-of-Experts Model Serving with Disaggregated Expert Parallelism (\textit{SIGCOMM 2025})
\item KVServe: Service-Aware KV Cache Compression for Communication-Efficient Disaggregated LLM Serving (\textit{SIGCOMM 2026})
}



\projectarea{Network Measurement \& Operations}{
    \item \lr{Telemetry \& Monitoring}
    \item \lr{Measurement \& Characterization}
    \item \lr{Security, Privacy \& Trust}
    \item \lr{Diagnosis \& Root Cause Analysis}
    \item \lr{Packet Processing \& Data Structures}
    \item \lr{Routing \& Path Selection}
    \item \lr{Programmability \& Hardware Offload}
    \item \lr{Control-Plane Design \& Management}
    \item \lr{Topology, Fabric \& Architecture}
    \item \lr{Traffic Engineering \& Load Balancing}
}{observability, telemetry, troubleshooting, configuration/state management, network digital twins, P4, SDN, SRv6}{
\item Raha: A General Tool to Analyze WAN Degradation (\textit{SIGCOMM 2025})
\item Raising the Level of Abstraction for Sketch-Based Network Telemetry with SketchPlan (\textit{IMC 2024})
}

\projectarea{Programmable Network Infrastructure}{
    \item \lr{Programmability \& Hardware Offload}
    \item \lr{Service Chaining \& Virtualization}
    \item \lr{Packet Processing \& Data Structures}
    \item \lr{Security, Privacy \& Trust}
    \item \lr{Scheduling \& Resource Allocation}
    \item \lr{Control-Plane Design \& Management}
    \item \lr{Placement, Offloading \& Partitioning}
    \item \lr{Measurement \& Characterization}
    \item \lr{Reliability, Fault Tolerance \& Recovery}
    \item \lr{Verification \& Correctness}
}{programmable data planes, NFV, service chaining, in-network computing, middleboxes, P4, eBPF, SmartNIC/DPU, NFV, service-function chaining, SDN}{
\item OptimusPrime: Unleash Dataplane Programmability through a Transformable Architecture (\textit{SIGCOMM 2024})
\item VeriNC: Finding Design Risks of In-Network Computing Systems (\textit{ICNP 2026})
}

\projectarea{Internet Routing, WAN \& CDN}{
    \item \lr{Routing \& Path Selection}
    \item \lr{Measurement \& Characterization}
    \item \lr{Security, Privacy \& Trust}
    \item \lr{Architecture \& System Design}
    \item \lr{Traffic Engineering \& Load Balancing}
    \item \lr{Caching \& Content Placement}
    \item \lr{Control-Plane Design \& Management}
    \item \lr{Coding, Compression \& Redundancy}
    \item \lr{Traffic Analysis \& Classification}
    \item \lr{Transport \& Congestion Control}
}{interdomain routing, WANs, naming/addressing, CDNs, BGP, DNS, IPv6, Anycast, submarine cables, Internet scanning, Tor/onion services}{
\item Threading the Ocean: Mapping Digital Routes Across Submarine Cables using Calypso (\textit{SIGCOMM 2025})
\item Who Holds the Steering Wheel? Opacity and Consolidation in CDN Replica Selection (\textit{NINeS 2026})
}

\projectarea{P2P \& Decentralized Networks}{
    \item \lr{Security, Privacy \& Trust}
    \item \lr{Distributed Coordination \& Consensus}
    \item \lr{Throughput, Capacity \& Scalability}
    \item \lr{Routing \& Path Selection}
    \item \lr{Scheduling \& Resource Allocation}
    \item \lr{Protocol Design \& Interoperability}
    \item \lr{Reliability, Fault Tolerance \& Recovery}
    \item \lr{Topology, Fabric \& Architecture}
    \item \lr{Coding, Compression \& Redundancy}
    \item \lr{Latency \& Deadline Optimization}
}{peer-to-peer overlays, blockchain networks, decentralized services, Blockchain/P2P, payment channels, anonymous communication}{
\item Deanonymizing Ethereum Validators: The P2P Network Has a Privacy Issue (\textit{USENIX Security 2025})
\item Reliable and Decentralized Certificate Revocation via DNS: The Case for RevDNS (\textit{SIGCOMM 2025})
}

\projectarea{Quantum Networks}{
    \item \lr{Routing \& Path Selection}
    \item \lr{Scheduling \& Resource Allocation}
    \item \lr{Topology, Fabric \& Architecture}
    \item \lr{Cost, Economics \& Provisioning}
    \item \lr{Throughput, Capacity \& Scalability}
    \item \lr{Control-Plane Design \& Management}
    \item \lr{Fairness, Isolation \& Multi-Tenancy}
    \item \lr{Measurement \& Characterization}
    \item \lr{Protocol Design \& Interoperability}
    \item \lr{Verification \& Correctness}
}{quantum Internet, entanglement distribution, quantum switches, distributed quantum systems, quantum entanglement}{
\item Optimal Scheduling in a Quantum Switch (\textit{SIGMETRICS 2025})
\item Quantum Network Optimization: From Optimal Routing to Fair Resource Allocation (\textit{SIGMETRICS 2025})
}
\noindent{\footnotesize\textbf{شرط انتخاب:} تنها پروژه‌های \textbf{شبکه‌محور}، مانند مسیریابی، زمان‌بندی، تخصیص منابع یا ارزیابی پروتکل‌های شبکه‌ی کوانتومی پذیرفته‌اند؛ پروژه‌های صرفاً محاسبات کوانتومی خارج از دامنه‌اند.}\par




\end{document}
